% Generated draft: requires mathematical and source review.
\documentclass[11pt]{article}
\usepackage[margin=25mm]{geometry}
\usepackage{fontspec}
\setmainfont{Latin Modern Roman}
\usepackage{amsmath,amssymb,mathtools,unicode-math}
\setmathfont{Latin Modern Math}
\usepackage{enumitem,needspace,xurl,hyperref}
\setlength{\parindent}{0pt}
\setlength{\parskip}{5pt}
\setlength{\emergencystretch}{3em}
\newcommand{\N}{\mathbb{N}}
\newcommand{\Z}{\mathbb{Z}}
\newcommand{\Q}{\mathbb{Q}}
\newcommand{\R}{\mathbb{R}}
\newcommand{\C}{\mathbb{C}}
\newcommand{\F}{\mathbb{F}}
\newcommand{\A}{\mathbb{A}}
\newcommand{\PP}{\mathbb{P}}
\newcommand{\E}{\mathbb{E}}
\newcommand{\eps}{\varepsilon}
\newcommand{\dd}{\,\mathrm{d}}
\begin{document}
\section*{UnsolvedMath ID: 10800008}
\textbf{Source problem code:} catalogue entry
\subsection{Definitions and mathematical statement}
Retain the source notation \(f_a,g_b\) and \(U_i,\varphi_i\). To state its approximation requirement, let \(P\) be the intended parameter space of pairs \((a,b)\), and let \((X,d)\) be the metric space in which the roots lie. Set
\[
 R(a,b)=\{x\in X:f_a(x)=0, g_b(x)=0\},\qquad
 N_\varepsilon(R(a,b))=\bigcup_{r\in R(a,b)}B_d(r,\varepsilon).
\]
The problem asks the same questions as in the missing preceding problem, with exact root selections replaced by approximate selections satisfying, for one fixed \(\varepsilon>0\),
\[
 \varphi_i(a,b)\in N_\varepsilon(R(a,b))
 \quad\text{for every }i\text{ and every }(a,b)\in U_i.
\]
The distance is to an actual common root, rather than a bound on the residuals \(f_a,g_b\).

The supplied summary describes the intended context as continuous selections on an open cover \(P=\bigcup_iU_i\) and asks about the approximate covering number, its relation to the exact number, and the possibility of a single global approximate selection. In that interpretation, for a fixed tolerance define the approximate number to be the smallest number of open sets admitting continuous maps with the displayed property, or \(+\infty\) if no finite such cover exists. Number \(1\) means a continuous approximate selection on all of \(P\). The actual preceding questions, the equations and their root/parameter domains are absent; this contextual formulation does not reconstruct them. The record also does not say whether the question is for one specified tolerance, some tolerance, every tolerance or all sufficiently small tolerances.
\subsection{Short English statement}
Replace the missing system's exact continuous root selections by points within one fixed positive distance of some common root, and ask the preceding covering/selection questions for this approximate notion.
\subsection{Sources}
\begin{itemize}
\item Victor Vassiliev, \emph{A Few Problems on Monodromy and Discriminants}, Problem 2B. \url{https://amj.math.stonybrook.edu/html-articles/Files-2015-2024/15-11/}. Recorded access 2026-07-29; extraction html-structure. The preceding system and questions are not included in the supplied extract.
\item Local UnsolvedMath registry, \texttt{data/unsolvedmath/problems.json}, permanent ID \(10800008\), exact source code \texttt{catalogue entry}. Its OPEN-TRIAGE summary describes exact versus approximate covering numbers and possible global approximate selection, without an additional resolving reference.
\end{itemize}
\end{document}
